\documentclass[10pt,a4paper]{amsart}
\usepackage[margin=19mm]{geometry}
\usepackage{amsmath,amssymb,mathtools}
\usepackage[scr=rsfs,cal=euler]{mathalfa}
\usepackage{xfrac}
\usepackage[unicode,hidelinks,bookmarks=false]{hyperref}
\newcommand{\ie}{i.e.\ }
\newcommand{\cf}{cf.\ }
\newcommand{\resp}{respectively\ }
\newcommand{\Subsection}{\subsection}
\providecommand{\nobreakdash}{\nobreak\hbox{-}}
\setlength{\emergencystretch}{2em}
\multlinegap0pt
\newtheorem{mainthm}{Theorem}
\newtheorem{prop}{Proposition}
\usepackage{cleveref}
\crefname{mainthm}{Theorem}{Theorems}
\crefname{prop}{Proposition}{Propositions}
\DeclareMathOperator{\Coh}{\catfont{Coh}}
\DeclareMathOperator{\Dmat}{\mathbb{D}^{\mathrm{Mat}}}
\DeclareMathOperator{\Dse}{\mathbb{D}^{\mathrm{Se}}}
\newcommand{\F}{\mathcal{F}}
\renewcommand{\O}{\mathcal{O}}
\DeclareMathOperator{\QCoh}{\catfont{QCoh}}
\renewcommand{\S}{\mathcal{S}}
\newcommand{\Y}{\mathcal{Y}}
\newcommand{\catfont}{\mathrm}
\DeclareMathOperator{\cone}{cone}
\newcommand{\localcoh}{\Gamma_{BG}}
\newcommand{\punct}[1]{{#1}^\circ}
\DeclareMathOperator{\uHom}{\underline{\mathcal{H}om}}
\title{Serre Functors and Local Duality for Affine Quotients}
\author{Ivan Noden}
\date{Source: arXiv:2605.21290v2 (CC BY 4.0)}
\begin{document}
\maketitle

\noindent\emph{Source and licence.} Serre Functors and Local Duality for Affine Quotients. Version arXiv:2605.21290v2 (CC BY 4.0), distributed by arXiv under the Creative Commons Attribution 4.0 International licence, http://creativecommons.org/licenses/by/4.0/, as recorded in the arXiv licence metadata for this exact version and shown on the abstract page https://arxiv.org/abs/2605.21290v2. Full licence notice: \texttt{licenses/arxiv-2605.21290-CC-BY-4.0.txt}.\par\smallskip

\noindent\textit{Editorial scope.} Counted unit: the article's prop prop:dmat\_coh\_se (source0.tex lines 1090--1092). The article's complete original proof of it is delivered with it (source0.tex lines 1093--1121, one \texttt{proof} environment of 29 lines). No mathematics is written, repaired or re-derived: the statement and the proof are the article's own lines, in the article's own notation, and no retained line is substituted or re-flowed.  Retained uncounted as the source-local context the proof needs: lines 210--212; lines 764--767; lines 1066--1068.  Nothing was omitted from any retained span, so the package carries no omission record.  No retained line contains a citation, so the wrapper carries no bibliography and no bibliography entry.  No retained line contains a figure environment, an image-inclusion command or drawing code, so no image is delivered and the package declares no asset dependency.  Editorial formatting only, in addition: an AMS \texttt{amsart} preamble that restates the article's own declarations and macros used by the retained lines, namely the environments \texttt{mainthm}, \texttt{prop} on separate counters, together with the standard packages those lines need.  The retained environments print in this document as Theorem 1, Proposition 1 in order of appearance, whereas the article prints the same items with the numbering of its own full document.  The complete native source of the article as distributed in the arXiv e-print for this exact version is bundled unchanged as \texttt{sources/201054/source0.tex}.
  \begin{mainthm}\label{thm:serre_functor}
    We have \(\S_\Y \simeq \localcoh \Psi(\omega_\Y)[\dim G]\).
  \end{mainthm}

  \begin{prop}\label{prop:k_generates}
    The functor \(\uHom_{\QCoh(\Y)}(i_\ast \O_{BG}, -)\) is conservative
    on \(\Coh(\Y)\).
  \end{prop}

  \begin{prop}\label{prop:dloc_coh}
    If \(\F \in \Coh(\Y)\), then \(\Dmat(\localcoh \F) \in \Coh(\Y)\).
  \end{prop}

  \begin{prop}\label{prop:dmat_coh_se}
    For all \(\F \in \Coh(\Y)\), there is a natural equivalence \(\Dmat(\localcoh(\F)) \simeq \Dse(\F)[\dim G]\).
  \end{prop}

  \begin{proof}
    Observe that 
    \begin{displaymath}
      \uHom_{\QCoh(\Y)}(\localcoh \F, \S_\Y) \simeq \uHom_{\QCoh(\Y)}(\F, \uHom_{\QCoh(\Y)}(\localcoh \O_\Y, \S_\Y))
    \end{displaymath}
    since \(\localcoh \F \simeq \localcoh \O_\Y \otimes \F\). Therefore, it
    is enough to show \(\Dmat(\localcoh \O_\Y) \simeq \Psi(\omega_\Y)[\dim G]\).
    From \Cref{thm:serre_functor}, we know \(\S_\Y \simeq \localcoh \Psi(\omega_\Y)[\dim G]\)
    and therefore, as \(\localcoh \O_\Y \in \QCoh(\Y)_{BG}\),
    \begin{displaymath}
      \uHom_{\QCoh(\Y)}(\localcoh \O_\Y, \S_\Y) \simeq \uHom_{\QCoh(\Y)}(\localcoh \O_\Y, \Psi(\omega_\Y)[\dim G]) \simeq \Dse(\localcoh \O_\Y)[\dim G].
    \end{displaymath}

    Now, note
    \begin{displaymath}
      \Dse(\localcoh \O_\Y) \simeq \cone(\Dse(j_\ast \O_{\punct{\Y}}) \to \Psi(\omega_\Y))
    \end{displaymath}
    so the claim is proven if we show \(\Dse (j_\ast \O_{\punct{\Y}})\simeq 0\).
    However, \(\Psi(\omega_\Y) \in \Coh(\Y)\) and \(\Dse(\localcoh \O_\Y) \simeq \Dmat(\localcoh \O_\Y)[-\dim G] \in \Coh(\Y)\)
    by \Cref{prop:dloc_coh} so we also know 
    \begin{displaymath}
      \Dse(j_\ast \O_{\punct{\Y}}) \in \Coh(\Y).
    \end{displaymath}
    We are then done since
    \begin{displaymath}
      \uHom_{\QCoh(\Y)}(i_\ast \O_{BG}, \Dse(j_\ast \O_{\punct{\Y}})) \simeq \uHom_{\QCoh(\Y)}(j_\ast j^\ast i_\ast \O_{BG}, \Psi(\omega_\Y)) \simeq 0
    \end{displaymath}
    by base change and this implies \(\Dse(j_\ast \O_{\punct{\Y}}) \simeq 0\) by \Cref{prop:k_generates}
  \end{proof}

\end{document}
